\section{Conclusiones y Trabajo Futuro}
\label{sec:concl}

\subsection{Contribuciones}

\begin{enumerate}
  \item \textbf{Arquitectura eficiente}: se propone una arquitectura de dos
        redes ($\approx$18K parámetros) con Fourier features e imposición
        estructural de $\partial \zb/\partial t = 0$, frente a los
        $\approx$544K de la arquitectura monolítica de Dazzi.
        \pendiente{cuantificación de su precisión frente a A2 y A3 en el
        estudio de escalamiento (Sección~\ref{sec:arch-comp}).}

  \item \textbf{Equifinalidad}: se caracterizan cuatro mecanismos para romper
        la degeneración $(h, \zb)$ del problema inverso ---observaciones de
        velocidad, muestreo temporal, condiciones de borde y la pérdida
        $\mathcal{L}_{\zb,\text{IC}}$---; su eficacia se evalúa en los casos
        de prueba (Sección~\ref{sec:exp}).

  \item \textbf{Formas de ecuación para inversión}: primer estudio comparativo
        de las formas primitiva, primitiva-conservativa y conservativa del
        residuo SWE en el problema inverso (Sección~\ref{sec:ablation}).
        \pendiente{conclusión sobre qué forma del residuo es más robusta para
        la inversión.}

  \item \textbf{Robustez}: el banco de pruebas evalúa la tolerancia del método
        al ruido en las observaciones y a la densidad de muestreo.
        \pendiente{niveles de ruido y densidad de observaciones tolerados.}

  \item \textbf{Límite con datos reales (Exp.\ 6)}: se evalúa de forma
        reproducible si el PINN de penalización blanda recupera la batimetría
        de los datos reales de Angel et al., frente al método adjunto de
        restricción dura. La barrera potencial es estructural ---la
        equifinalidad $\eta = h + \zb$ bajo física blanda con sensores fuera
        del bump---, no computacional.
        \pendiente{recuperación del PINN sobre los datos de Angel y
        comparación con el benchmark adjunto.}
\end{enumerate}

\subsection{Trabajo futuro}

\begin{itemize}
  \item \textbf{Restricción dura de la EDP}: motivado por la dificultad del
        PINN de penalización blanda sobre datos reales (Exp.\ 6), integrar un
        solver directo SWE diferenciable en el lazo (híbrido PINN--adjunto),
        entrenamiento causal/curriculum, o ponderación adaptativa de residuos,
        para cerrar la brecha con el benchmark adjunto de Angel sobre datos
        reales.
  \item \textbf{Formas de ecuación}: extender la comparación de formas del
        residuo al caso transitorio 2D con wetting-drying.
  \item \textbf{Co-inversión de fricción}: aprender simultáneamente $\zb$
        y el coeficiente de Manning $n$.
  \item \textbf{Escalabilidad}: evaluar factibilidad en dominios 2D reales
        (río Sacramento, $\sim$10K puntos).
\end{itemize}
